% !TeX spellcheck = en_US
\documentclass[12pt]{article}

\usepackage{times,fullpage,xspace,fancyhdr,url,color}
\usepackage[pdftex]{graphicx}
\usepackage[pdftex,
            colorlinks=true,
            urlcolor=black,
            linkcolor=black,
            citecolor=black,
            bookmarksopen=false,
            bookmarksnumbered=true,
            pdfstartview=FitH]{hyperref}

\input{common/variables}
\input{common/dates}

% Fitness-test-specific fallback values if not supplied by shared variables.
\providecommand{\MaxTimePassingLeaderboard}{120}

\pdfcompresslevel=9
\hypersetup{
 pdftitle={\leaguename The Fitness Tests},
 pdfauthor={Technical Committee HSL},
}
\usepackage{microtype}
\usepackage[utf8]{inputenc}
\usepackage{amsmath}
\usepackage{xargs}
\usepackage[colorinlistoftodos,prependcaption,textsize=tiny]{todonotes}
\usepackage{siunitx}
\usepackage[capitalize,noabbrev]{cleveref}
\usepackage[official]{eurosym}
\usepackage[useregional]{datetime2}
\usepackage{subcaption}
\usepackage{enumitem}
\usepackage{xcolor}
\usepackage{tikz}
\usetikzlibrary{calc,math}
\DTMlangsetup[en-GB]{ord=raise,monthyearsep={,\space}}

\newcommandx{\unsure}[2][1=]{\todo[linecolor=red,backgroundcolor=red!25,bordercolor=red,#1]{#2}}
\newcommandx{\change}[2][1=]{\todo[linecolor=blue,backgroundcolor=blue!25,bordercolor=blue,#1]{#2}}
\newcommandx{\info}[2][1=]{\todo[linecolor=green,backgroundcolor=green!25,bordercolor=green,#1]{#2}}
\newcommandx{\improvement}[2][1=]{\todo[linecolor=Plum,backgroundcolor=Plum!25,bordercolor=Plum,#1]{#2}}

% comment 'disable' in to disable all the todo notes :)
\usepackage
[
%disable
]{todonotes}

\usepackage[theorems]{tcolorbox}
\newtcbtheorem[number within=section]{hintbox}{}%
{colback=red!10,colframe=red!45!black,fonttitle=\bfseries}{th}

\sloppy
\newcommand{\ie}{\mbox{i.\,e.}\xspace}
\newcommand{\eg}{\mbox{e.\,g.}\xspace}
%\newcommand{\cf}{\mbox{cf.}\xspace}
\newcommand{\cf}{see\xspace}
% \newcommand{\comment}[1]{\marginpar{\pdfannot width 4in height .5in depth 8pt {/Subtype /Text /Contents (#1)}}}
\newcommand{\inparagraph}[1]{\paragraph{#1\hspace{-1em} }}


% some colors
\definecolor{orange}{rgb}{1,0.5,0}
\definecolor{red}{rgb}{1,0,0}
\definecolor{green}{rgb}{0,1,0}


\title{The \leaguename\\ Fitness Tests}
\author{RoboCup Technical Committee}
\date{(\RCYear fitness tests, as of \today)}

\setlength{\parindent}{0pt}
\setlength{\parskip}{12pt plus 6pt minus 3 pt}
\setcounter{tocdepth}{1}
\widowpenalty=10000
\clubpenalty=10000

\pagestyle{fancy}
\lhead{}
\chead{}
\rhead{}
\lfoot{}
\cfoot{}
\rfoot{}

\renewcommand{\headrulewidth}{0.4pt}
\renewcommand{\footrulewidth}{0.4pt}

% needed to align an image and text correctly side by side
\newcommand{\imagebox}[1]{\raisebox{2ex}{\raisebox{-\height}{#1}}}

% Shared field base comes from the same macros used by the main rules document.
% Shared TikZ helpers for rulebook and fitness test field visualizations.

\providecommand{\FieldAnnotation}[5]{%
  \draw[thick] ($#1$) -- ($#1 + (#3:#4)$);
  \draw[thick] ($#2$) -- ($#2 + (#3:#4)$);
  \draw[<->, thick] ($#1 + (#3:#4/2)$) -- ($#2 + (#3:#4/2)$);
  \node at ($#1!0.5!#2 + (#3:#4)$) {#5};
}

\providecommand{\FieldDrawGoal}[4]{%
  \draw[goalLine]
  ($#1 + (90+#4:#2/2)$) --
  ($#1 + (90+#4:#2/2) + (180+#4:#3+\goalLineWidth/2)$) --
  ($#1 + (-90+#4:#2/2) + (180+#4:#3+\goalLineWidth/2)$) --
  ($#1 + (-90+#4:#2/2)$);
}

\providecommand{\FieldDrawPenaltyMark}[2]{%
  \fill[white] #1 circle (#2/2);
}

\providecommand{\SetFieldCanvas}[2]{%
  \def\fieldCanvasHalfLengthValue{#1}%
  \def\fieldCanvasHalfWidthValue{#2}%
}

\providecommand{\SetFieldGeometryScale}[1]{%
  \def\fieldGeometryScaleValue{#1}%
}

\providecommand{\SetFieldDrawingScale}[1]{%
  \def\fieldDrawingScaleValue{#1}%
}

\providecommand{\ConfigureSmallField}{%
  \def\fieldDrawingScaleValue{0.0125}%
  \def\fieldLengthValue{9.0}%
  \def\fieldWidthValue{6.0}%
  \def\goalDepthValue{0.5}%
  \def\goalWidthValue{1.8}%
  \def\goalAreaLengthValue{1.0}%
  \def\goalAreaWidthValue{3.0}%
  \def\penaltyAreaLengthValue{2.0}%
  \def\penaltyAreaWidthValue{4.0}%
  \def\penaltyMarkDistValue{1.5}%
  \def\penaltyMarkLengthValue{0.10}%
  \def\cornerArcRadiusValue{0.0}%
  \def\centerCircleDiameterValue{1.5}%
  \def\fieldLineWidthValue{0.05}%
  \def\goalLineWidthValue{0.10}%
  \def\borderStripValue{1.0}%
  \SetFieldGeometryScale{1.0}%
  \SetFieldCanvas{5.5}{4.0}%
}

\providecommand{\ConfigureMediumField}{%
  \def\fieldDrawingScaleValue{0.01}%
  \def\fieldLengthValue{14.0}%
  \def\fieldWidthValue{9.0}%
  \def\goalDepthValue{0.7}%
  \def\goalWidthValue{2.4}%
  \def\goalAreaLengthValue{1.0}%
  \def\goalAreaWidthValue{4.0}%
  \def\penaltyAreaLengthValue{3.0}%
  \def\penaltyAreaWidthValue{6.0}%
  \def\penaltyMarkDistValue{2.0}%
  \def\penaltyMarkLengthValue{0.10}%
  \def\cornerArcRadiusValue{0.5}%
  \def\centerCircleDiameterValue{3.0}%
  \def\fieldLineWidthValue{0.05}%
  \def\goalLineWidthValue{0.10}%
  \def\borderStripValue{1.0}%
  \SetFieldGeometryScale{1.0}%
  \SetFieldCanvas{8.0}{5.5}%
}

\providecommand{\ConfigureLargeField}{%
  \def\fieldDrawingScaleValue{0.006666}%
  \def\fieldLengthValue{22.0}%
  \def\fieldWidthValue{14.0}%
  \def\goalDepthValue{0.6}%
  \def\goalWidthValue{2.4}%
  \def\goalAreaLengthValue{1.0}%
  \def\goalAreaWidthValue{4.0}%
  \def\penaltyAreaLengthValue{3.5}%
  \def\penaltyAreaWidthValue{7.0}%
  \def\penaltyMarkDistValue{2.5}%
  \def\penaltyMarkLengthValue{0.15}%
  \def\cornerArcRadiusValue{1.0}%
  \def\centerCircleDiameterValue{4.0}%
  \def\fieldLineWidthValue{0.12}%
  \def\goalLineWidthValue{0.20}%
  \def\borderStripValue{1.0}%
  \SetFieldGeometryScale{1.0}%
  \SetFieldCanvas{12.0}{8.0}%
}

\providecommand{\DrawConfiguredField}[1]{%
  \begin{tikzpicture}[
    x=\fieldDrawingScaleValue*100cm,
    y=\fieldDrawingScaleValue*100cm,
    fieldLine/.style = {line width=\fieldDrawingScaleValue*5cm, color=white},
    goalLine/.style = {line width=\fieldDrawingScaleValue*10cm, color=white}
  ]
    \tikzmath{
      \borderStrip=\borderStripValue;
      \fieldLength=\fieldLengthValue;
      \fieldWidth=\fieldWidthValue;
      \goalDepth=\goalDepthValue;
      \goalWidth=\goalWidthValue;
      \goalAreaLength=\goalAreaLengthValue;
      \goalAreaWidth=\goalAreaWidthValue;
      \penaltyAreaLength=\penaltyAreaLengthValue;
      \penaltyAreaWidth=\penaltyAreaWidthValue;
      \penaltyMarkDist=\penaltyMarkDistValue;
      \penaltyMarkLength=\penaltyMarkLengthValue;
      \cornerArcRadius=\cornerArcRadiusValue;
      \centerCircleDiameter=\centerCircleDiameterValue;
      \fieldLineWidth=\fieldLineWidthValue;
      \goalLineWidth=\goalLineWidthValue;
      \arenaWidth=\fieldWidth + 2 * \borderStrip;
      \arenaLength=\fieldLength + 2 * \borderStrip;
      \hflw=\fieldLineWidth/2;
    }

    \path[use as bounding box]
      (-\fieldCanvasHalfLengthValue,-\fieldCanvasHalfWidthValue)
      rectangle
      (\fieldCanvasHalfLengthValue,\fieldCanvasHalfWidthValue);

    \begin{scope}[scale=\fieldGeometryScaleValue]
      \coordinate (center) at (0,0);
      \coordinate (fieldBL) at ($(-\fieldLength/2,-\fieldWidth/2)$);
      \coordinate (fieldBR) at ($(\fieldLength/2,-\fieldWidth/2)$);
      \coordinate (fieldTL) at ($(-\fieldLength/2,\fieldWidth/2)$);
      \coordinate (fieldTR) at ($(\fieldLength/2,\fieldWidth/2)$);
      \coordinate (fieldLineBL) at ($(fieldBL) + (\hflw,\hflw)$);
      \coordinate (fieldLineBR) at ($(fieldBR) + (-\hflw,\hflw)$);
      \coordinate (fieldLineTL) at ($(fieldTL) + (\hflw,-\hflw)$);
      \coordinate (fieldLineTR) at ($(fieldTR) + (-\hflw,-\hflw)$);
      \coordinate (goalL) at ($(-\fieldLength/2,0)$);
      \coordinate (goalR) at ($(\fieldLength/2,0)$);
      \coordinate (goalLineL) at ($(goalL) + (\hflw,0)$);
      \coordinate (goalLineR) at ($(goalR) + (-\hflw,0)$);
      \coordinate (penaltyMarkL) at ($(-\fieldLength/2 + \penaltyMarkDist,0)$);
      \coordinate (penaltyMarkR) at ($(\fieldLength/2 - \penaltyMarkDist,0)$);
      \coordinate (goalAreaLineL) at ($(-\fieldLength/2 + \goalAreaLength - \hflw,0)$);
      \coordinate (goalAreaLineR) at ($(\fieldLength/2 - \goalAreaLength + \hflw,0)$);

      \fill[black!40!green] ($(-\arenaLength/2,-\arenaWidth/2)$) rectangle ($(\arenaLength/2,\arenaWidth/2)$);

      \draw[fieldLine] ($(fieldBL) + (\hflw,\hflw)$) rectangle ($(fieldTR) - (\hflw,\hflw)$);

      \ifdim \cornerArcRadius pt > 0pt
        \draw[fieldLine] (fieldBL) + (0:\cornerArcRadius) arc (0:90:\cornerArcRadius);
        \draw[fieldLine] (fieldBR) + (90:\cornerArcRadius) arc (90:180:\cornerArcRadius);
        \draw[fieldLine] (fieldTR) + (180:\cornerArcRadius) arc (180:270:\cornerArcRadius);
        \draw[fieldLine] (fieldTL) + (270:\cornerArcRadius) arc (270:360:\cornerArcRadius);
      \fi

      \draw[fieldLine] ($(0,-\fieldWidth/2 + \hflw)$) rectangle ($(0,\fieldWidth/2 - \hflw)$);
      \FieldDrawGoal{(goalL)}{\goalWidth}{\goalDepth}{0};
      \FieldDrawGoal{(goalR)}{\goalWidth}{\goalDepth}{180};

      \draw[fieldLine] ($(-\fieldLength/2 + \hflw,-\goalAreaWidth/2 + \hflw)$)
        rectangle
        ($(-\fieldLength/2 - \hflw + \goalAreaLength,\goalAreaWidth/2 - \hflw)$);
      \draw[fieldLine] ($(\fieldLength/2 - \hflw,-\goalAreaWidth/2 + \hflw)$)
        rectangle
        ($(\fieldLength/2 + \hflw - \goalAreaLength,\goalAreaWidth/2 - \hflw)$);

      \draw[fieldLine] ($(-\fieldLength/2 + \hflw,-\penaltyAreaWidth/2 + \hflw)$)
        rectangle
        ($(-\fieldLength/2 - \hflw + \penaltyAreaLength,\penaltyAreaWidth/2 - \hflw)$);
      \draw[fieldLine] ($(\fieldLength/2 - \hflw,-\penaltyAreaWidth/2 + \hflw)$)
        rectangle
        ($(\fieldLength/2 + \hflw - \penaltyAreaLength,\penaltyAreaWidth/2 - \hflw)$);

      \FieldDrawPenaltyMark{(penaltyMarkL)}{\penaltyMarkLength}
      \FieldDrawPenaltyMark{(penaltyMarkR)}{\penaltyMarkLength}
      \draw[fieldLine] (center) circle (\centerCircleDiameter/2 - \fieldLineWidth/2);
      \FieldDrawPenaltyMark{(center)}{\penaltyMarkLength}

      \coordinate (goalLFront) at ($(goalL) + (0:\fieldLineWidth)$);
      \coordinate (goalLBack) at ($(goalL) + (180:\goalDepth)$);
      \coordinate (goalLBottom) at ($(goalL) - (90:\goalWidth/2 - \goalLineWidth/2)$);
      \coordinate (goalLTop) at ($(goalL) + (90:\goalWidth/2 - \goalLineWidth/2)$);
      \coordinate (goalAreaLT1) at ($(-\fieldLength/2,\goalAreaWidth/2)$);
      \coordinate (goalAreaLT2) at ($(goalAreaLT1) + (0:\goalAreaLength)$);
      \coordinate (goalAreaLT3) at ($(goalAreaLT2) + (-90:\goalAreaWidth)$);
      \coordinate (goalAreaRT1) at ($(\fieldLength/2,\goalAreaWidth/2)$);
      \coordinate (goalAreaRT2) at ($(goalAreaRT1) + (180:\goalAreaLength)$);
      \coordinate (goalAreaRT3) at ($(goalAreaRT2) + (-90:\goalAreaWidth)$);
      \coordinate (penaltyAreaLT1) at ($(\fieldLength/2,\penaltyAreaWidth/2)$);
      \coordinate (penaltyAreaLT2) at ($(penaltyAreaLT1) + (180:\penaltyAreaLength)$);
      \coordinate (penaltyAreaLT3) at ($(penaltyAreaLT2) + (-90:\penaltyAreaWidth)$);
      \coordinate (centerA1) at ($(center) + (-45:\centerCircleDiameter/2)$);
      \coordinate (centerA2) at ($(center) + (135:\centerCircleDiameter/2)$);
      \coordinate (borderStripA1) at ($(fieldTL) + (90:\borderStrip)$);
      \coordinate (borderStripA2) at ($(fieldTL) + (180:\borderStrip)$);
      \coordinate (cornerArcMeasureStart) at ($(fieldTR) + (180:\cornerArcRadius)$);

      #1
    \end{scope}
  \end{tikzpicture}%
}

\definecolor{leaderactive}{rgb}{0,0,1}
\definecolor{leadertarget}{rgb}{0,0.75,0}
\definecolor{leaderpath}{rgb}{1,0,1}
\definecolor{leaderball}{rgb}{1,0.9,0}

\newcommand{\FitnessTestDot}[2]{\node[circle, fill=#2, draw=black, line width=1.0pt, inner sep=0pt, minimum size=11pt, transform shape=false] at #1 {};}
\newcommand{\FitnessTestLabel}[3]{\node[font=\Large\bfseries, text=#3, transform shape=false] at #1 {#2};}
\newcommand{\FitnessTestArrow}[2]{\draw[-{Latex[length=11pt,width=8pt]}, color=leaderpath, line width=2.0pt] #1 -- #2;}

\newcommand{\FitnessTestFieldS}[1]{%
  \begingroup
  \ConfigureSmallField
  \SetFieldDrawingScale{0.01}
  \SetFieldGeometryScale{2.444444}
  \SetFieldCanvas{13.6}{10.0}
  \DrawConfiguredField{#1}%
  \endgroup
}

\newcommand{\FitnessTestFieldM}[1]{%
  \begingroup
  \ConfigureMediumField
  \SetFieldDrawingScale{0.01}
  \SetFieldGeometryScale{1.571429}
  \SetFieldCanvas{13.6}{10.0}
  \DrawConfiguredField{#1}%
  \endgroup
}

\newcommand{\FitnessTestFieldL}[1]{%
  \begingroup
  \ConfigureLargeField
  \SetFieldDrawingScale{0.01}
  \SetFieldGeometryScale{1.0}
  \SetFieldCanvas{13.6}{10.0}
  \DrawConfiguredField{#1}%
  \endgroup
}


\begin{document}

\maketitle

\begin{center}
Questions or comments on the fitness test rules should be submitted via \url{https://github.com/RoboCup-HumanoidSoccerLeague/HSL-Rules/issues}, to the \texttt{\#rule-book} channel on the RoboCup Official Discord server, or by mail to \url{rc-hsl-tc@lists.robocup.org}.
\end{center}

\newpage

\tableofcontents
\setcounter{tocdepth}{3}

\thispagestyle{fancy}

\clearpage

\cfoot{\thepage}
\setcounter{page}{1}

\section{Introduction}
Starting in 2026, a fitness test will be introduced to track and compare team performances across key skills over multiple years. 
These fitness tests track the development of the league and are used to make informed decisions on rule changes such as field sizes or game duration.
They further aim to motivate teams to improve in specific areas valuable to humanoid robotic soccer and to compile a record of the top-performing components and skills. 
By highlighting these individual skills, the league encourages focused development in areas that contribute to overall team performance, and creates a showcase of the best technical abilities in the league.

\textbf{The \RCYear fitness tests will be a trial run to ensure feasibility, and as such will not have any impact on team scores and will not award any certificates or trophies. Failing to complete these fitness tests will not have any negative repercussions. However, all teams are required to show up to their assigned slots and must declare if they are unable to complete any of the challenges.}

\subsection{Challenge Rules}

The fitness tests are conducted as part of the robot inspection and consist of three individual challenges as described below. Each challenge can be attempted up to three times and should be attempted at least once. If a robot is unable to complete any of the challenges, the team may declare this and thus skip the challenge. A skipped challenge is marked as "not attempted". If a challenge is attempted multiple times, the best result is counted.

Challenges require a single robot. If a team has multiple different robot platforms, they may ask for permission to make a second attempt at the fitness test with another robot platform, but are not required to do so. The OC may reject this request if the schedule does not allow for it. Teams are not allowed to use multiple different robot platforms in the same fitness test attempt.

\subsection{Setup}

The fitness tests are conducted on a small division field. Additional required materials include a stopwatch, a tape measure and one ball of each division.

\subsection{Results}

The results of the fitness tests will be recorded and published on the Humanoid Soccer League website (hsl.robocup.org). Results will be kept indefinitely to allow analysis of the development of the league over time.

Unless otherwise specified, all results are recorded in SI base units with 2 decimals of accuracy.

% !TeX root = ../HSL-Fitness-Tests.tex
% !TeX spellcheck = en_US

\section{Fastest Walk Test}
This fitness test measures the fastest straight autonomous walk from one side of the field to the other.

\subsection{Setup}
This fitness test takes place on one half of a standard SPL field and requires 1 active robot.

The competing robot is placed 50cm behind the starting touchline in line with the penalty mark.
The run ends when the robot crosses the opposite touchline.

\begin{figure}[ht]
    \centering
    \centering
\resizebox{\linewidth}{!}{%
  \FitnessTestFieldM{%
    \draw[orange, line width=4.0pt] (fieldLineTL) -- (fieldLineTR);
    \FitnessTestDot{($(penaltyMarkR) + (0,-\fieldWidth/2 - 0.5)$)}{leaderactive}
  }%
}

    \caption{Fastest Walk Fitness Test (M-Field): the start point is shown in blue and the finishing touchline is highlighted.}
    \label{fig:walk_test_m}
\end{figure}

\subsection{Challenge Execution}
In each attempt, the active robot must walk autonomously in a straight run from the starting side of the field to the opposite touchline.
Code changes are not allowed during an attempt, but are allowed between attempts.

\subsection{Scoring}
The head referee starts the timer once the robot touches the starting touchline and stops the timer once the target touchline is touched.

The lowest time across all attempts is recorded, rounded to the nearest second.

\section{Kick Accuracy Fitness Test}
This fitness test measures the accuracy with which a robot can kick an HSL competition soccer ball toward specified targets.

\subsection{Setup}
This fitness test takes place on a standard SPL field and requires 1 active robot and 1 ball.
The ball is placed on the center of the edge of the goal area.
The competing robot is placed in the center of the goal line facing the ball.

\begin{figure}[ht]
    \centering
    \centering
\resizebox{\linewidth}{!}{%
  \FitnessTestFieldM{%
    \FitnessTestDot{(goalLineR)}{leaderactive}
    \FitnessTestDot{(goalAreaLineR)}{leaderball}
    \FitnessTestDot{(center)}{leadertarget}
    \FitnessTestLabel{($(center) + (-0.45,-0.45)$)}{A}{black}
    \FitnessTestDot{(fieldLineTL)}{leadertarget}
    \FitnessTestLabel{($(fieldLineTL) + (0.55,-0.55)$)}{B}{black}
    \FitnessTestDot{(fieldLineBL)}{leadertarget}
    \FitnessTestLabel{($(fieldLineBL) + (0.55,0.55)$)}{C}{black}
  }%
}

    \caption{Kick Accuracy Fitness Test (M-Field): robot is shown in blue, targets are shown in green, and ball placement is shown in yellow.}
    \label{fig:kick_accuracy_test_m}
\end{figure}

\subsection{Challenge Execution}
Each attempt consists of kicking toward each of the three targets (illustrated in \cref{fig:kick_accuracy_test_m}; S-Field and L-Field variants are provided in \cref{fig:kick_accuracy_test_s,fig:kick_accuracy_test_l}):
\begin{itemize}
    \item The center mark in the center circle
    \item The left corner on the opposite side of the field
    \item The right corner on the opposite side of the field
\end{itemize}
In each attempt, for each kick, participating teams must start the robot behavior of their competing robot via a single button press. The robot then has \qty{\PenaltyShootoutKickTime}{\second} to kick the ball towards the current target.
The competing robot must not touch the ball a second time after the ball has clearly moved, otherwise the attempt ends immediately without scoring.
Once the ball has come to a complete stop, the distance from the ball to the target is measured.

Code changes are allowed when changing targets.

\subsection{Scoring}
After all attempts, the best distance achieved for each target is recorded. 
Accuracy is scored inversely based on distance: the closer the ball stops to the target, the higher the score.
Each target is scored independently, with the lowest distance measurement across attempts recorded for each target.
Distances are rounded to the closest half-centimetre.

\section{Recovery from Fall Fitness Test}
This fitness test measures how quickly a robot can recover and resume normal operation after falling forward or backward.

\subsection{Setup}
This fitness test takes place on a standard SPL field and requires 1 active robot.
The robot is placed on the ground in the center of the field.
Each run starts from a defined ground pose: one run from a front-fall pose and one run from a back-fall pose.

\subsection{Challenge Execution}
Each attempt starts with the robot already on the ground in the assigned pose (flat on the front or back).
The robot must then autonomously recover and stand up.
The test is performed with two separate attempts: one forward fall and one backward fall.

For each fall direction, teams start the robot recovery behavior via a single button press.
Code changes between forward and backward fall attempts are allowed.

\subsection{Scoring}
The head referee starts timing at the button press, and stops the timer when the robot is standing stably on both feet.
The fastest recovery time across all attempts for each fall direction is recorded, rounded to the nearest tenth of a second.
The scores for forward and backward recovery are summed to produce the final fitness test score.

\appendix

\section{Field Size Variants}

\subsection{Fastest Walk Test Variants}
The S-Field and L-Field variants of the Fastest Walk Test are shown in \cref{fig:walk_test_s,fig:walk_test_l}.

\begin{figure}[ht]
    \centering
    \centering
\resizebox{\linewidth}{!}{%
  \FitnessTestFieldS{%
    \draw[orange, line width=4.0pt] (fieldLineTL) -- (fieldLineTR);
    \FitnessTestDot{($(penaltyMarkR) + (0,-\fieldWidth/2 - 0.5)$)}{leaderactive}
  }%
}

    \caption{Fastest Walk Fitness Test (S-Field): the start point is shown in blue and the finishing touchline is highlighted.}
    \label{fig:walk_test_s}
\end{figure}

\begin{figure}[ht]
    \centering
    \centering
\resizebox{\linewidth}{!}{%
  \FitnessTestFieldL{%
    \draw[orange, line width=4.0pt] (fieldLineTL) -- (fieldLineBL);
    \FitnessTestDot{(($(penaltyAreaLT1) + (0.5,0)$)}{leaderactive}
  }%
}

    \caption{Fastest Walk Fitness Test (L-Field): the start point is shown in blue and the finishing touchline is highlighted.}
    \label{fig:walk_test_l}
\end{figure}

\subsection{Kick Accuracy Fitness Test Variants}
The S-Field and L-Field variants of the Kick Accuracy Fitness Test are shown in \cref{fig:kick_accuracy_test_s,fig:kick_accuracy_test_l}.

\begin{figure}[ht]
    \centering
    \input{figs/fitness_tests/kick_accuracy_test_s}
    \caption{Kick Accuracy Fitness Test (S-Field): robot is shown in blue, targets are shown in green, and ball placement is shown in yellow.}
    \label{fig:kick_accuracy_test_s}
\end{figure}

\begin{figure}[ht]
    \centering
    \input{figs/fitness_tests/kick_accuracy_test_l}
    \caption{Kick Accuracy Fitness Test (L-Field): robot is shown in blue, targets are shown in green, and ball placement is shown in yellow.}
    \label{fig:kick_accuracy_test_l}
\end{figure}


\end{document}
